\chapter{自演化数字孪生}
\label{ch:evolve}

本章要判定的量是：一次过程结束时，什么允许写回下一次的先验，以及在线预测偏离物理降阶到什么程度就必须丢弃修正项。演化发生在参数和残差项上，不发生在能量方程的结构上。

\section{离线算子与在线更新的分工}

\paragraph{要判定的量}
在线一步预测里，哪一项是物理降阶，哪一项是有界残差。

\paragraph{继承的文献层与不继承的层}
神经算子作为参数化演化的加速器，继承 \cite{li2021fno,lu2021deeponet}。物理残差的约束继承物理信息网络里“方程不能被任意拟合”的要求 \cite{raissi2019}，但在线不把残差网络当成方程的发现者。不继承某一次训练的测试误差，因为没有这次训练。

\begin{equation}
\label{eq:op}
\vect{z}_{k+1}=\mathcal{G}_{\phi}(\vect{z}_k,\vect{u}_k,\vect{\theta})+\vect{r}_{\psi}(\vect{z}_k,\vect{u}_k).
\end{equation}
\(\mathcal{G}_{\phi}\) 在仿真轨迹上训练。Fourier 神经算子与 DeepONet 使用同一训练集。主模型取蒸汽干预位移误差更低者，另一套保留为对照。二者都未训练，所以没有主模型。

残差约束
\begin{equation}
\label{eq:bnp}
\|\vect{r}_{\psi}\|\le r_{\max}(\vect{z},\vect{u}).
\end{equation}
\(r_{\max}\) 由离线验证集的误差分位数确定。超出该界的输出被截断，并回退到两节点或一维物理步进。这是“一次异常过程不能改写能量结构”的机制。

本次在线预测就是十六条控制体本身，没有 \(\mathcal{G}_{\phi}\)，也没有 \(\vect{r}_{\psi}\)。回退条件因此不会被触发，也不能据此声称物理约束已经受过检验。

\paragraph{本节不能推出的结论}
不能比较两种算子的位移误差，不能声明胜者。

\section{贝叶斯更新}

\paragraph{要判定的量}
下一次同类别的先验是这一次后验的覆盖，还是精度加权。

类别先验由离线仿真和公开热物性范围构成。窗内似然来自 \eqref{eq:obs}。
\begin{equation}
\label{eq:bayes}
p(\vect{\theta}\mid\vect{y}_{0:w},\mathrm{class})
\propto p(\vect{y}_{0:w}\mid\vect{\theta})\,p(\vect{\theta}\mid\mathrm{class}).
\end{equation}
过程结束时，把均值、协方差、类别、厚度区间和探针是否有效写回本地知识库。下一次的先验是精度加权，不是用一条轨迹覆盖类别：
\begin{equation}
\label{eq:hier}
\begin{aligned}
\vect{\Lambda}_{\mathrm{next}}&=\vect{\Lambda}_{\mathrm{hyper}}+\sum_{j\in\mathrm{class}}\vect{\Lambda}_{j},\\
\vect{\Lambda}_{\mathrm{next}}\vect{\mu}_{\mathrm{next}}
&=\vect{\Lambda}_{\mathrm{hyper}}\vect{\mu}_{\mathrm{hyper}}+\sum_{j}\vect{\Lambda}_{j}\vect{\mu}_{j}.
\end{aligned}
\end{equation}
后验协方差比先验更大时，本次不写回。类别识别置信度低时，使用各类别中最保守的进度分位数，并提高探针倾向。

本次战役结束于日志文件，没有知识库，没有 \(\vect{\Lambda}\) 的更新。19440 行不能累加成一个类别后验：它们是同一格式的因子样本，而且热阻更新使用了真值参照。

\paragraph{本节不能推出的结论}
精度公式不说明厨房之间的数据可以合并。知识库若存在，也只是本地的。本次连本地写回都没有。

\section{一次烹饪内部的演化}

\paragraph{要判定的量}
内部热阻在几十分钟内的漂移，是否比类别差异更快，因而必须在过程内部跟踪。

规格
\begin{equation}
R_{\mathrm{int},k+1}=R_{\mathrm{int},k}+\kappa_c\langle M_c-\varepsilon_{l,s}\rangle_+\Delta t+\xi_k
\end{equation}
与 \eqref{eq:structure} 一致。探针在漂移后可以再触发一次，间隔不短于 \(t_{\mathrm{gap}}\)，避免探针本身成为主要加热方式。

结壳增量只在表面含水低于 0.45 时发生。S1 的初始含水不低于 0.64，表面失水率约 \(10^{-11}\,\mathrm{s}^{-1}\)，结壳停在 0。导热场本来也不读这个变量。热阻代理只在 190 s 附近更新一次，没有第二次探针，也没有 \(t_{\mathrm{gap}}\) 的实现。种子 7 的方差在 192 s 变为 0.224 之后保持到时域终点。过程后段没有新的结壳信息可以再进入估计。

\paragraph{本节不能推出的结论}
不能把“一次更新”写成热阻漂移已经被在线跟踪。漂移的规格需要多次更新。

\section{算子训练数据如何覆盖干预}

\paragraph{要判定的量}
训练集里有没有蒸汽阶跃。若只有恒温跟踪，算子在干预上会外推失败，因果替换输入就不可信。

规格协议：

\begin{enumerate}
\item 食材为肉糜、土豆、面团。
\item 几何为厚度三档、初温两档、含水两档。这里的含水档少于 S1 的三档，不能把两套格子当成同一个设计。
\item 执行器为加热、蒸汽、风速、排汽的拉丁超立方，并插入 P1--P4。
\item 每条轨迹保存真值决策状态、带噪观测和执行器。
\item 划分按厚度分组，禁止同一厚度的相邻轨迹同时出现在训练和干预测试中。
\end{enumerate}

教师轨迹的规格规模是高保真约 360 条，降阶扩充到 \(10^4\) 量级。扩充条目标记 \texttt{augmented=true}，不进入主对照的分母。本次日志全部是 \texttt{augmented=false}，但也全部不是教师轨迹：它们是四条固定策略，不是拉丁超立方，也不含算子的训练标签。S0 的网格减半、S2 的失效注入、S3 的教师集都没有作为独立套件运行。

\paragraph{本节不能推出的结论}
S1 的 4860 个组合不能拿来训练算子后再用同一批做干预测试。那样会把内插误差写成外推能力。

\section{失效与回退}

三条硬限制：物理残差在滚动窗口上超过 \(r_{\max}\) 时丢弃残差项；后验精度不高于先验时不写回；类别置信度低时采用最保守的分位数。降级的方向只允许朝向更少依赖估计器，详见第 \ref{ch:safe} 章。

本次没有残差项，所以第一条没有对象。没有写回，所以第二条没有对象。类别由实验设计给定，不存在识别置信度。回退次数因此不能画成数据图。图位保留坐标，不绘制曲线。

\paragraph{本节不能推出的结论}
“本次没有越界”不是“越界率为零”的实验结果。没有监测项就不能计算比率。
